\documentclass[12pt]{article}

\usepackage[margin=1in]{geometry}
\usepackage[T1]{fontenc}
\usepackage[utf8]{inputenc}
\usepackage{booktabs}
\usepackage{tabularx}
\usepackage{array}
\usepackage{longtable}
\usepackage{rotating}
\usepackage{pdflscape}
\usepackage{hyperref}
\usepackage{enumitem}
\usepackage{caption}

\hypersetup{
  bookmarks=true,
  colorlinks=true,
  linkcolor=blue,
  citecolor=green,
  filecolor=magenta,
  urlcolor=cyan
}

\setlength{\parindent}{0pt}
\setlength{\parskip}{0.75em}
\renewcommand{\arraystretch}{1.15}

\newcolumntype{Y}{>{\raggedright\arraybackslash}X}

\title{Hazard Analysis\\Project Proxi}
\author{%
  Team \#23\\
  Savinay Chhabra\\
  Amanbeer Singh Minhas\\
  Gourob Podder\\
  Ajay Singh Grewal
}
\date{April 6, 2026}

\begin{document}

\maketitle
\thispagestyle{empty}
\newpage

\pagenumbering{roman}

\begin{table}[h!]
\centering
\caption{Revision History}
\label{tab:revhist}
\begin{tabularx}{\textwidth}{l l Y}
\toprule
\textbf{Date} & \textbf{Ver.} & \textbf{Notes} \\
\midrule
10/10/2025 & 1.0 & Added initial hazard analysis. \\
06/04/2026 & 1.1 & Reformatted revision history to one item per row. \\
06/04/2026 & 1.1 & Added in-text references to all tables. \\
06/04/2026 & 1.1 & Replaced ruled tables with Booktabs styling. \\
06/04/2026 & 1.1 & Rotated the FMEA table for readability. \\
06/04/2026 & 1.1 & Added external LLM privacy hazard and mitigations. \\
06/04/2026 & 1.1 & Aligned hazard requirement IDs with the SRS and V\&V \\
           &     & terminology. \\
\bottomrule
\end{tabularx}
\end{table}

\newpage
\tableofcontents
\listoftables
\listoffigures
\newpage

\pagenumbering{arabic}

\section{Introduction}

In Project Proxi, a hazard is any system condition that can lead to harm,
loss, or an unacceptable outcome during normal use, foreseeable misuse, or
failure. For this project, the most relevant harms are not physical injury,
but unintended file changes, privacy loss, accessibility failure, incorrect
automation, and loss of user trust.

The purpose of this document is to identify important hazards early and to
connect them to mitigations and requirements that remain consistent with the
Software Requirements Specification (SRS), Module Guide (MG), Module
Interface Specification (MIS), and Verification and Validation (V\&V)
artifacts. Table~\ref{tab:revhist} summarizes the evolution of this
revision.

The rest of the document is organized as follows. Section~2 defines the
scope and purpose of the analysis. Section~3 describes system boundaries and
major components. Section~4 records critical assumptions. Section~5 presents
our Failure Modes and Effects Analysis (FMEA) in
Table~\ref{tab:fmea}. Section~6 derives safety and security requirements.
Finally, Section~7 prioritizes the mitigation roadmap using
Tables~\ref{tab:ship} and~\ref{tab:later}.

\section{Scope and Purpose of Hazard Analysis}

Proxi is an AI-powered desktop assistant that supports voice and text input
for accessible, high-level computer interaction. Its baseline scope includes
speech input, text input, natural-language interpretation, task planning,
controlled execution through MCP tools or desktop automation, and textual or
spoken feedback.

This hazard analysis focuses on software and interaction hazards within that
scope. It covers failures in speech and text handling, intent
interpretation, confirmation logic, automation, logging, privacy handling,
external service use, and accessibility support. It also considers hazards
introduced by interactions with external applications and services when they
are part of a supported Proxi workflow.

The following items are outside the direct scope of this analysis:
\begin{itemize}[noitemsep]
  \item physical hardware failure of the user's computer,
  \item failures of the operating system itself,
  \item outages of unrelated third-party applications,
  \item failures of the wider network infrastructure, and
  \item specialized medical or life-critical device control.
\end{itemize}

Even with these exclusions, the consequences of software hazards are still
significant. Incorrect actions may delete or alter files, disclose personal
information, reduce accessibility for target users, or make the system feel
unsafe to use. The goal of this document is therefore to identify credible
hazards, estimate their effects, and derive practical controls that fit the
actual Proxi architecture.

\section{System Boundaries and Components}

\subsection{System Boundaries}

Proxi operates at the level of high-level desktop assistance. Inside the
system boundary are the components that accept user input, interpret intent,
plan actions, enforce confirmation and permission checks, execute supported
actions, and present results.

Within the system boundary, Proxi:
\begin{itemize}[noitemsep]
  \item accepts input through speech and text,
  \item processes requests using intent interpretation and task planning,
  \item invokes MCP tools or desktop automation for supported actions,
  \item stores limited local state such as logs and settings, and
  \item returns visible or spoken feedback to the user.
\end{itemize}

Outside the system boundary are the operating system, physical hardware,
third-party applications, and cloud services that Proxi may depend on but
does not fully control. This means hazards can arise both from internal
logic and from interfaces with adjacent systems.

\subsection{Major System Components}

For hazard analysis purposes, the most important conceptual components are:
\begin{enumerate}[label=\arabic*., noitemsep]
  \item \textbf{User Interface Module:} accepts speech or text input and
  presents feedback.
  \item \textbf{Natural Language Understanding Engine:} converts raw input
  into structured intent.
  \item \textbf{Task Planning and Execution Manager:} decides how to carry
  out the requested task.
  \item \textbf{Tool and System Integration Layer:} invokes MCP tools and
  desktop automation.
  \item \textbf{Context and Session Manager:} keeps short-term interaction
  state.
  \item \textbf{Data Storage and Logging Component:} stores logs, settings,
  and support information.
  \item \textbf{Security and Permissions Module:} enforces confirmation,
  permission scope, and access policies.
  \item \textbf{External AI Service Interface:} sends approved task-relevant
  data to external LLM or speech services when those services are enabled.
\end{enumerate}

The last component is especially important in this revision. Because the SRS
allows the use of off-the-shelf LLM services, privacy hazards must consider
what information is transmitted externally and whether that transmission is
necessary and consented to.

\section{Critical Assumptions}

This analysis is based on the following assumptions:
\begin{enumerate}[label=A\arabic*., noitemsep]
  \item Proxi runs on supported desktop systems such as Windows and macOS
  using standard consumer hardware.
  \item Users have a working microphone and display, and can fall back to
  text interaction when voice input is unreliable.
  \item Supported MCP tools, automation adapters, and adjacent applications
  generally behave according to their documented interfaces.
  \item Internet connectivity may be unavailable or unstable, but local
  workflows still remain possible where designed.
  \item High-risk actions are expected to pass through explicit confirmation
  logic before execution.
  \item When external LLM or speech services are used, Proxi sends only the
  minimum task-relevant content needed for the requested operation, not a
  continuous full-screen stream by default.
  \item Users are informed when a workflow requires external processing of
  personal or sensitive content.
\end{enumerate}

Assumptions A6 and A7 were made explicit in this revision because privacy
risk depends strongly on whether Proxi transmits only selected content or an
excessive amount of desktop context.

\section{Failure Mode and Effects Analysis}

Table~\ref{tab:fmea} presents the main FMEA results. The table is rotated to
improve readability and to avoid overly compressed text.

\clearpage
\begin{landscape}
\small
\setlength{\LTleft}{0pt}
\setlength{\LTright}{0pt}

\begin{longtable}{p{2.7cm} p{2.9cm} p{3.0cm} c p{3.2cm} c c p{4.0cm}
  p{2.2cm} p{1.5cm}}
\caption{Failure Modes and Effects Analysis (FMEA)}
\label{tab:fmea}\\
\toprule
\textbf{Component} &
\textbf{Failure Mode} &
\textbf{Effect of Failure} &
\textbf{S} &
\textbf{Cause of Failure} &
\textbf{O} &
\textbf{D} &
\textbf{Recommended Action} &
\textbf{SRS Ref.} &
\textbf{Hazard ID} \\
\midrule
\endfirsthead

\multicolumn{10}{l}{\textbf{Table \thetable{} continued from previous
page}}\\
\toprule
\textbf{Component} &
\textbf{Failure Mode} &
\textbf{Effect of Failure} &
\textbf{S} &
\textbf{Cause of Failure} &
\textbf{O} &
\textbf{D} &
\textbf{Recommended Action} &
\textbf{SRS Ref.} &
\textbf{Hazard ID} \\
\midrule
\endhead

\midrule
\multicolumn{10}{r}{\textbf{Continued on next page}}\\
\endfoot

\bottomrule
\endlastfoot

Speech-to-Text (STT) &
Misinterprets voice commands &
Wrong action or no action taken &
8 &
Background noise, accent variation, or poor microphone quality &
5 &
6 &
Improve STT robustness, provide transcript visibility, and allow
correction before execution &
POA.R.1 &
H1-1 \\

Command Interpreter &
Fails to understand user intent &
Task not completed or incorrect action &
7 &
Ambiguous phrasing or missing context &
4 &
6 &
Add clarification prompts and stronger context handling &
FUNC.R.3 &
H2-1 \\

Tool / MCP Integration &
Wrong tool is triggered &
Unintended application or function runs &
8 &
Incorrect tool mapping, stale registry entry, or version mismatch &
4 &
6 &
Validate tool selection, constrain tool scope, and provide safer
fallbacks &
SAF.R.2, IAS.R.1 &
H3-1 \\

Desktop Automation &
Command crashes, freezes, or affects the wrong target &
User task fails or unrelated application state changes &
7 &
Permission errors, fragile selectors, or unresponsive applications &
5 &
5 &
Add timeouts, retry logic, targeted application checks, and graceful
failure handling &
ROFT.R.2, WER.R.1 &
H4-1 \\

Confirmation Layer &
High-risk action executes without approval &
Accidental deletion, overwrite, or external action &
9 &
Logic bug, missing risk tag, or skipped validation path &
3 &
7 &
Require explicit confirmation for privileged or destructive actions &
SAF.R.1, FUNC.R.9 &
H5-1 \\

Accessibility Layer &
Captions, screen reader flow, or other assistive support fails &
Users with disabilities cannot complete supported tasks &
9 &
Missing WCAG support, poor semantics, or inaccessible feedback timing &
4 &
6 &
Test with assistive technologies and validate against WCAG 2.2 and
AODA expectations &
ACC.R.2 &
H6-1 \\

Feedback Manager &
Response is delayed, missing, or unclear &
User confusion, repeated commands, or loss of trust &
5 &
High load, TTS failure, or poor status messaging &
5 &
5 &
Provide visible status indicators and clear, plain-language responses &
SAL.R.1, UAP.R.1 &
H7-1 \\

Update and Deployment &
New release breaks a working feature &
Supported commands stop working or become inconsistent &
6 &
Incomplete regression testing or packaging errors &
4 &
5 &
Use regression testing, rollback support, and versioned release
artifacts &
LON.R.2, REL.R.1 &
H8-1 \\

Logging and Storage &
Logs are missing, corrupted, or over-collected &
Traceability is lost or unnecessary user data is retained &
6 &
Logging disabled, storage failure, or excessive capture
configuration &
4 &
5 &
Keep audit logs reliable, bounded, and privacy-aware &
AUD.R.1, PRIV.R.2 &
H9-1 \\

Cross-platform Compatibility &
Behaviour differs across operating systems &
Features fail on some supported platforms &
7 &
API differences, permission models, or UI variation across OSes &
4 &
5 &
Test on supported Windows and macOS targets and isolate
platform-specific adapters &
EPE.R.1, ADP.R.1 &
H10-1 \\

External AI Service Use &
More information than necessary is sent to an external LLM or speech
service &
Privacy loss, policy non-compliance, or user distrust &
9 &
Poor prompt construction, over-broad context capture, or missing
consent &
3 &
7 &
Minimize outbound data, require user consent, and send only the
minimum task-relevant text, file excerpt, or structured fields
needed for the request &
PRIV.R.2, IAS.R.2 &
H11-1 \\

External AI Service Use &
Sensitive content is sent from the wrong source, such as an unintended
window or document &
Disclosure of confidential data unrelated to the task &
9 &
Incorrect context selection, stale window focus, or unsafe
automation binding &
2 &
8 &
Require explicit target selection for sensitive operations and avoid
whole-screen capture by default &
PRIV.R.2, SAF.R.2 &
H11-2 \\

\end{longtable}

\vspace{0.3cm}
\noindent\textbf{Legend:} S = Severity (1--10), O = Occurrence (1--10),
D = Detection (1--10).

\end{landscape}
\clearpage

The highest-risk items in Table~\ref{tab:fmea} are the failure to confirm
high-risk actions and the privacy hazards introduced by external AI service
use. These hazards received the highest severity because they can cause data
loss or disclosure even when the rest of the system appears to be operating
normally.

The new privacy rows clarify the TA's concern directly. In the intended
Proxi design, the system should not send the user's whole screen by default.
Instead, it should send only the minimum information needed for the specific
task, such as the current command text, a selected file excerpt, or a small,
explicitly chosen content segment.

\section{Derived Safety and Security Requirements}

The hazards in Table~\ref{tab:fmea} lead to the following requirements.
These are phrased to remain consistent with the terminology and intent of the
SRS, MG, MIS, and V\&V plan.

\subsection{Safety Requirements}

\begin{description}[style=nextline]
  \item[SAF.R.1 - Explicit confirmation for risky actions]
  Any privileged, destructive, or externally impactful action shall require
  explicit user confirmation before execution.

  \textit{Why:} This mitigates H5-1 and reduces the chance of silent harmful
  execution.

  \textit{Verification idea:} Test cases for deletes, overwrites, sends, and
  removals should show a confirmation step every time.

  \item[SAF.R.2 - Scoped and validated execution]
  Proxi shall validate targets, tool choice, and execution scope before
  running automation or external integrations.

  \textit{Why:} This mitigates H3-1, H4-1, and H11-2.

  \textit{Verification idea:} Out-of-scope targets, risky destinations, and
  ambiguous application targets should trigger blocking or clarification.

  \item[ROFT.R.2 - Graceful failure handling]
  Supported workflows shall fail safely when tools, applications, or network
  dependencies are unavailable.

  \textit{Why:} This mitigates H4-1 and H8-1.

  \textit{Verification idea:} Simulated timeouts and unavailable adapters
  should end in a clear failure state rather than partial unsafe execution.

  \item[ACC.R.2 - Accessible recovery and feedback]
  Safety prompts, status messages, and recovery paths shall remain accessible
  to users who rely on voice, captions, scaling, or assistive technologies.

  \textit{Why:} A safety control that cannot be perceived or understood does
  not adequately control risk for the target users.

  \textit{Verification idea:} Confirmation and error flows should be checked
  with screen readers, captions, and enlarged text.
\end{description}

\subsection{Security and Privacy Requirements}

\begin{description}[style=nextline]
  \item[PRIV.R.2 - Data minimization for external services]
  When Proxi uses an external LLM, STT, or TTS service, it shall send only
  the minimum task-relevant data needed to complete the request.

  \textit{Why:} This directly mitigates H11-1 and aligns with the SRS rule
  that only strictly needed data should leave the user's computer.

  \textit{Verification idea:} Inspect representative outbound requests and
  confirm that they contain only the selected text, command, or excerpt
  rather than broad desktop context.

  \item[PRIV.R.3 - Explicit consent for external sharing]
  Proxi shall clearly inform the user when a workflow requires data to be
  processed by an external AI service and shall obtain consent before sending
  personal or sensitive content.

  \textit{Why:} This mitigates H11-1 and supports informed user control.

  \textit{Verification idea:} Privacy-sensitive workflows should present a
  consent step before transmission.

  \item[ACS.R.2 - Permission-bounded access]
  The system shall operate only within the permission scope explicitly
  granted by the user for folders, applications, and external actions.

  \textit{Why:} This mitigates H3-1, H4-1, and H11-2.

  \textit{Verification idea:} Requests outside allowed scope should be
  blocked or require a new permission grant.

  \item[AUD.R.1 - Traceable and privacy-aware audit logging]
  The system shall keep audit logs for actions, confirmations, and outcomes
  without silently storing unnecessary personal content.

  \textit{Why:} This mitigates H9-1 while preserving debugging and
  accountability.

  \textit{Verification idea:} Logs should record action metadata and outcome
  state, but should avoid excessive raw content retention.

  \item[IAS.R.2 - Protected communication with adjacent systems]
  External integrations shall use supported protected communication
  mechanisms, such as HTTPS, when network transport is required.

  \textit{Why:} This reduces exposure to interception and tampering during
  external service use.

  \textit{Verification idea:} Network inspection should confirm encrypted
  transport for supported cloud-backed workflows.
\end{description}

\section{Roadmap}

Table~\ref{tab:ship} prioritizes the mitigation work that most directly
reduces risk in the capstone scope. Table~\ref{tab:later} lists useful
follow-on work that is valuable but less critical for the final milestone.

\begin{table}[h!]
\centering
\caption{Highest-Priority Mitigations for the Capstone Scope}
\label{tab:ship}
\begin{tabularx}{\textwidth}{l Y l}
\toprule
\textbf{ID} & \textbf{Mitigation to Implement} & \textbf{Target Window} \\
\midrule
SAF.R.1 & Confirmation step for destructive, privileged, or externally
impactful actions & Sprint 1 \\
SAF.R.2 & Stronger target validation and scoped execution checks & Sprint 1 \\
PRIV.R.2 & Data minimization for outbound LLM and speech-service requests &
Sprint 1--2 \\
PRIV.R.3 & Consent flow for privacy-sensitive external processing &
Sprint 2 \\
AUD.R.1 & Reliable and bounded audit logging for actions and outcomes &
Sprint 2 \\
ROFT.R.2 & Safe timeout, retry, and failure handling in adapters & Sprint 2 \\
\bottomrule
\end{tabularx}
\end{table}

\begin{table}[h!]
\centering
\caption{Deferred or Longer-Term Mitigations}
\label{tab:later}
\begin{tabularx}{\textwidth}{l Y Y}
\toprule
\textbf{ID} & \textbf{Mitigation} & \textbf{Reason for Deferral} \\
\midrule
ACS.R.2 & Deeper sandboxing and finer-grained permission profiles & Requires
more platform-specific engineering than fits the capstone timeline \\
IAS.R.2 & Broader automated inspection of all external service traffic & Best
added after the baseline supported workflows are stabilized \\
REL.R.1 & More comprehensive release hardening and signing workflow & Useful
for long-term deployment, but not strictly required for the course release \\
\bottomrule
\end{tabularx}
\end{table}

\clearpage

\section*{Appendix - Reflection Questions}
\addcontentsline{toc}{section}{Appendix - Reflection Questions}

The purpose of these reflection questions is to help the team assess what it
learned while creating the hazard analysis and how the process could improve
in future projects.

\begin{enumerate}
  \item What went well while writing this deliverable?
  \item What pain points did you experience during this deliverable, and how
  did you resolve them?
  \item Which risks had your team already considered before this deliverable,
  and which ones emerged during the hazard analysis itself?
  \item Other than physical harm, what other categories of software risk are
  important to consider, and why?
\end{enumerate}

\section{Reflections}

\subsection{Amanbeer Minhas Reflection}

\begin{enumerate}
  \item What went well while writing this deliverable?\\
  Writing this deliverable went well because we already had a clear idea of
  our system from the SRS and FMEA work. Breaking the project into
  components helped me think about hazards more systematically and kept the
  analysis organized. I also felt more confident in identifying potential
  risks and writing about them in a structured way.

  \item What pain points did you experience during this deliverable, and how
  did you resolve them?\\
  One challenge was figuring out how detailed to make each section without
  overcomplicating the document. I also found it tricky to write about
  hazards without repeating the same points as the SRS. I resolved this by
  focusing on specific examples from our project, like MCP integration and
  OS automation, and connecting them directly to potential risks.

  \item Which of your listed risks had your team thought of before this
  deliverable, and which did you think of while doing this deliverable?\\
  We had already thought about risks like incorrect command execution and
  permission issues earlier in the project. However, risks such as operating
  system compatibility problems and deployment-related failures came up
  while writing this hazard analysis. They became more obvious once we broke
  the system into components and looked deeper into how each part could
  fail.

  \item Other than the risk of physical harm, list at least two other types
  of risk in software products. Why are they important to consider?\\
  One major risk is data loss or corruption, which can harm user trust and
  make the system unreliable. Another risk is accessibility failure, where
  users with limitations might not be able to use the software as intended.
  Both are important because they affect user confidence, usability, and the
  overall success of the product.
\end{enumerate}

\subsection{Gourob Podder Reflection}

\begin{enumerate}
  \item What went well while writing this deliverable?\\
  Identifying system boundaries and major components went smoothly because
  the project architecture was already clear from prior SRS work. Defining
  modules like the User Interface, Natural Language Understanding, and Agent
  Manager helped organize potential hazards systematically. Team
  collaboration also allowed us to quickly align on how risks could manifest
  in both software and user interactions.

  \item What pain points did you experience during this deliverable, and how
  did you resolve them?\\
  Determining the appropriate level of detail for the hazard analysis
  without drifting into system design was challenging. It was sometimes
  difficult to distinguish between a hazard and a general usability concern.
  We resolved this by using safety engineering definitions and focusing on
  sources of harm or adverse outcomes rather than design flaws. Abstract
  risks such as privacy or misinterpretation of user intent were addressed
  using scenario-based brainstorming.

  \item Which of your listed risks had your team thought of before this
  deliverable, and which did you think of while doing this deliverable?\\
  We had already recognized usability and privacy risks, like
  misinterpreted commands or unauthorized data access. New risks emerged
  during the hazard analysis, especially related to component interactions,
  such as potential failures in the Task Planning module or dependency on
  unreliable external APIs. These became apparent as we examined each
  subsystem and its interface boundaries.

  \item Other than the risk of physical harm, list at least two other types
  of risk in software products. Why are they important to consider?\\
  Privacy and data security risks, such as improper handling of sensitive
  user data, can harm user trust and violate legal requirements.
  Reliability and availability risks, like system downtime or malfunctions,
  can prevent users from completing essential tasks, especially for
  assistive technologies. Both are critical for maintaining usability,
  accessibility, and overall user confidence.
\end{enumerate}

\subsection{Savinay Chhabra Reflection}

\begin{enumerate}
  \item What went well while writing this deliverable?\\
  The provided template helped structure our hazards. Collaboration and
  work distribution went more smoothly, and we could relate hazards
  directly to the well-structured SRS document.

  \item What pain points did you experience during this deliverable, and how
  did you resolve them?\\
  Staying solution-neutral while analyzing hazards was difficult. The line
  between requirements and design is blurred, and we sometimes discussed
  solutions prematurely. We resolved this by refocusing on hazards rather
  than proposed solutions.

  \item Which of your listed risks had your team thought of before this
  deliverable, and which did you think of while doing this deliverable?\\
  Initially, we focused on privacy and legal risks associated with data
  breaches. While performing hazard analysis, we discovered additional
  potential risks and mitigation strategies that were not previously
  considered.

  \item Other than the risk of physical harm, list at least two other types
  of risk in software products. Why are they important to consider?\\
  Legal risks due to data breaches or violations of user privacy, and
  compliance risks from accessibility failures in design, are both
  important because they affect user trust, legal standing, and the
  usability of the product.
\end{enumerate}

\subsection{Ajay Singh Grewal Reflection}

\begin{enumerate}
  \item What went well while writing this deliverable?\\
  Turning the hazards into short, testable requirements worked well. By
  adding criteria under each requirement, it became easier to understand how
  to verify them later. The FMEA table also helped organize the hazards
  clearly, and breaking the system into components made it easier to think
  about specific risks.

  \item What pain points did you experience during this deliverable, and how
  did you resolve them?\\
  Choosing criteria that were realistically achievable within our time frame
  was tricky. We had to balance being thorough with being practical. We
  resolved this by focusing on the highest-impact hazards that were most
  relevant to our project scope and user needs.

  \item Which of your listed risks had your team thought of before this
  deliverable, and which did you think of while doing this deliverable?\\
  We had already discussed risks like confirmation for destructive actions
  and permission issues. While writing the hazard analysis, we identified
  additional risks like OS compatibility and deployment-related failures.

  \item Other than the risk of physical harm, list at least two other types
  of risk in software products. Why are they important to consider?\\
  One important risk is sending more data than needed, or sending it without
  user consent, because that breaks user trust and may discourage continued
  use. Another important risk is reliability failure, where a bad update or
  broken tool makes the system unusable. Both matter because they affect
  user trust, usability, and long-term adoption.
\end{enumerate}

\end{document}
